\scanpage{10}

donc $E'(M,N)$ est un complexe dont les invariants $\underline{H}^{i}$ sont ceux de $E(M,N)$ en degré $i+2$, et qui en degré $2$ donne lieu : une suite exacte
\[
0\longrightarrow\underline{\mathrm{Ext}}^{2}(M,N)\longrightarrow
\overbrace{\underline{H}^{2}\bigl(E'(M,N)\bigr)}^{P(M,N)}
\xrightarrow{\sigma}\underline{\mathrm{Hom}}(M,{}_{2}N)\longrightarrow 0
\tag*{\makebox[0pt][r]{\makebox[\displaywidth][l]{$(*)$}}}
\]

Heuristiquement, $E'(M,\unsure{N|\deleted{M}N|\deleted{N}N|\mathbb{N}})$ est le complexe qui exprime le ``2-champ de Picard \deleted{\unsure{``|s}}strict'' formé des 1-champs de Picard (pas nécessairement stricts) \inserted{``\unsure{épi-glés|épinglés}'' par $M,N$} sur des objets variables de $X$, en admettant que ta théorie pour les 1-champs de Picard \inserted{stricts} s'étend aux 2-champs de Picard stricts (ce qui pour moi ne fait guère de doute) ; de même $E(M,N)$ correspond aux champs de Picard str\uline{ict}s épi-glés par $M,N$.

La suite exacte $(*)$ se construit en tous cas canoniquement ``à la main''\unsure{,|)} où $\deleted{\unsure{\mathrm{l}|\mathrm{d}}\overset{\unsure{\cdot|\phantom{\cdot}}}{\mathrm{e}}}$ terme médian est le faisceau des \unsure{classes|choses} \inserted{à} \deleted{\unsure{d'|à}} ``équivalence'' près des

% archive sheet number: 9
